\documentclass[12pt]{article}

\usepackage{amsmath}
\usepackage{amssymb}
\usepackage{geometry}
\usepackage{booktabs}

\geometry{margin=1in}

\title{Lecture 4 Scribe \\ Joint Probability and Conditional Probability}
\author{}
\date{}

\begin{document}

\maketitle

\section{Introduction to Probability Theory}

\subsection{Experiment, Outcome, Event, Sample Space}

\textbf{Definition: Experiment}

An experiment is a procedure that produces a result.

Example: Tossing a coin five times $(E_5)$

Reasoning: Probability theory studies uncertain processes, so we define the process as an experiment.

\bigskip

\textbf{Definition: Outcome}

An outcome is a possible result of an experiment.

\[
\xi_1 = HHTHT
\]

Reasoning: Probability will be assigned to outcomes, so all possible results must be listed.

\bigskip

\textbf{Definition: Event}

An event is a set of outcomes.

Example:

Event $C$ = all outcomes with even number of heads

Reasoning: We are interested in conditions, not only single outcomes.

\bigskip

\textbf{Definition: Sample Space}

Sample space = set of all possible outcomes

\[
S = \{H, T\}
\]

Properties:

\begin{itemize}
\item Mutually exclusive
\item Collectively exhaustive
\end{itemize}

Reasoning: To assign probability, we must know all outcomes.

\bigskip

Types of sample space:

\begin{center}
\begin{tabular}{ll}
\toprule
Experiment & Type \\
\midrule
Coin flip & Discrete \\
Die roll & Discrete \\
Flip until tails & Countably infinite \\
Random number in [0,1] & Continuous \\
\bottomrule
\end{tabular}
\end{center}

Reasoning: Different spaces require different probability models.

\section{Axioms of Probability}

\textbf{Axiom 1}

\[
0 \le Pr(A) \le 1
\]

Reasoning: Probability cannot be negative or greater than 1.

\bigskip

\textbf{Axiom 2}

\[
Pr(S) = 1
\]

Reasoning: One outcome must occur.

\bigskip

\textbf{Axiom 3}

If $A \cap B = \emptyset$

\[
Pr(A \cup B) = Pr(A) + Pr(B)
\]

Reasoning: No overlap, so add probabilities.

\bigskip

For infinite events:

\[
Pr\left(\bigcup A_i\right) = \sum Pr(A_i)
\]

Reasoning: Total probability is sum of disjoint events.

\section{Corollary}

For finite events:

\[
Pr\left(\bigcup_{i=1}^{M} A_i\right)
=
\sum_{i=1}^{M} Pr(A_i)
\]

Reasoning: Finite case follows from axiom.

\section{Propositions}

\textbf{Complement}

\[
Pr(A^c) = 1 - Pr(A)
\]

Reasoning: A and complement form whole sample space.

\bigskip

\textbf{Subset}

If $A \subset B$

\[
Pr(A) \le Pr(B)
\]

Reasoning: B contains A.

\bigskip

\textbf{Union rule}

\[
Pr(A \cup B)
=
Pr(A) + Pr(B) - Pr(A \cap B)
\]

Reasoning: Intersection counted twice.

\section{Assigning Probability}

\subsection{Classical}

All outcomes equally likely.

Coin:

\[
Pr(H)=1/2
\]

Dice even:

\[
Pr = 3/6
\]

Two dice sum = 5

\[
Pr = 4/36
\]

Reasoning: Count favorable / total.

\subsection{Relative Frequency}

\[
Pr(A) =
\lim_{n \to \infty}
\frac{n_A}{n}
\]

Reasoning: Probability is long-run frequency.

\section{Joint Probability}

\[
Pr(A,B) = Pr(A \cap B)
\]

Reasoning: Both events occur.

\bigskip

Relative frequency:

\[
Pr(A,B)
=
\lim
\frac{n_{AB}}{n}
\]

Reasoning: Count both events.

\subsection{Card Example}

\[
Pr(A)=26/52
\]

\[
Pr(B)=40/52
\]

\[
Pr(C)=13/52
\]

\[
Pr(A,B)=20/52
\]

\[
Pr(A,C)=13/52
\]

\[
Pr(B,C)=10/52
\]

Reasoning: Count cards satisfying both.

\subsection{Costume Example}

\[
Pr = \frac{1}{4} \times \frac{4}{6} = \frac{1}{6}
\]

Reasoning: Independent selections.

\section{Conditional Probability}

\[
Pr(A|B)
=
\frac{Pr(A,B)}{Pr(B)}
\]

Reasoning: Sample space becomes B.

\subsection{Product Rule}

\[
Pr(A,B)
=
Pr(A|B)Pr(B)
=
Pr(B|A)Pr(A)
\]

\subsection{Three events}

\[
Pr(A,B,C)
=
Pr(C|A,B)
Pr(B|A)
Pr(A)
\]

Reasoning: Multiply step by step.

\section{Cards Without Replacement}

\[
Pr(B|A) = 12/51
\]

Reasoning: One card removed.

\section{Poker Flush}

\[
\frac{13}{52}
\times
\frac{12}{51}
\times
\frac{11}{50}
\times
\frac{10}{49}
\times
\frac{9}{48}
\]

Any flush:

\[
4 \times P
\]

Reasoning: Four suits.

\section{Missing Key Example}

\[
P(R|L^c)
=
\frac{P(R)}{1-P(L)}
=
\frac{0.4}{0.6}
=
\frac{2}{3}
\]

Reasoning: Right implies not left.

\section{Exam Summary}

Remember:

\begin{itemize}
\item Experiment
\item Outcome
\item Event
\item Sample space
\item Axioms
\item Complement rule
\item Union rule
\item Classical probability
\item Relative frequency
\item Joint probability
\item Conditional probability
\item Product rule
\item Chain rule
\item Card example
\item Dice example
\item Poker example
\item Missing key example
\end{itemize}

\end{document}